\chapter*{Introducció}
\addcontentsline{toc}{chapter}{Introducció}

\section*{Motivació i justificació del tema}
% Text de la motivació...
Aquest treball de recerca parteix de la necessitat de poder crear una alternativa més sostenible per reduir el consum de combustible en els vaixells dins del sector del comerç marítim arreu del món, 
ja que aquest representa vora un 90\% del transport de mercaderies globalment. 
D'aquesta manera, que si aconseguíssim fer aquest canvi per rebaixar el combustible utilitzat, suposaria un estalvi immens. 
A més, com que el combustible és finit, no podrem dependre sempre d'aquesta font d'energia. 
Per aquest motiu, aquest treball pretén mostrar una alternativa sostenible i reutilitzable.
\section*{Objectius de la recerca}
% Text dels objectius...
Un dels objectius de fer aquest projecte és poder analitzar com funcionen els diferents sistemes de propulsió marítima per als vaixells 
i comparar-los per veure quins són més adequats en cada context.
\section*{Hipòtesi}
% Text de la hipòtesi...
La hipòtesi que es planteja en aquest projecte, es preveu que els resultats d'analitzar 
aquestes tecnologies comparades amb altres sistemes de propulsió presentin un millor rendiment
en termes d'eficiència i cost econòmic a llarg termini. 
Respecte a la maqueta, es preveu que representarà correctament el funcionament 
de la propulsió, però que segurament es veurà afectada l'eficiència 
a falta de condicions reals.
\section*{Metodologia}
% Text de la metodologia...
